% =============================================================================
%  §II.3 — Correspondances : domaine, image, produit
%  Fragment inclus par main.tex (\input). Blocs subsection seulement.
% =============================================================================

\subsection{§II.3 --- Correspondances : domaine, image, produit (E~II.10--11)}
Cluster de cinq résultats, conditionnés (le cas échéant) par l'hypothèse honnête
$\mathrm{est\_graphe}(G):=(\forall z)(z\in G\Rightarrow z\text{ couple})$, dans
\texttt{.../ii\_3\_correspondances/}.
\begin{itemize}
  \item $G\subset \mathrm{pr}_1G\times\mathrm{pr}_2G$ (tout ensemble de couples est partie d'un produit) ;
  \item $G\langle\mathrm{pr}_1G\rangle=\mathrm{pr}_2G$ (image du domaine $=$ ensemble des valeurs), \emph{inconditionnel} ;
  \item $\mathrm{pr}_1G=\varnothing\ (\text{resp. }\mathrm{pr}_2G=\varnothing)\Rightarrow G=\varnothing$ ;
  \item corollaire $A\supset\mathrm{pr}_1G\Rightarrow G\langle A\rangle=\mathrm{pr}_2G$ ;
  \item involution $(G^{-1})^{-1}=G$ (le champ $\mathrm{est\_graphe}$ est load-bearing
        pour $G\subset(G^{-1})^{-1}$).
\end{itemize}
\textbf{Preuves (\Nstar).} Algèbre de couples (\texttt{couple\_reciproque},
\texttt{couple\_dans\_dom/img}, \texttt{couple\_dans\_produit}), images
(\texttt{\_inst\_image}, \texttt{image\_dans\_img}, \texttt{image\_croissante}),
extensionnalité $A1$. Domaine/image-ensemble : $\mathrm{pr}_1G=\mathrm{dom}(G)$,
$\mathrm{pr}_2G=\mathrm{img}(G)$ (axiomes \texttt{DOM}/\texttt{IMG}).

\subsection{§II.3 --- Composition : monotonie ; diagonale et identité (E~II.13)}
\textbf{Monotonie de la composée (Remarque, E~II.13).}
$\{\,G_1\subset G_2,\ G_1'\subset G_2'\,\}\vdash G_1'\circ G_1\subset G_2'\circ G_2$
(\texttt{composee\_monotone}). Un couple $(p,r)\in G_1'\circ G_1$ provient d'un
$y$ avec $(p,y)\in G_1$, $(y,r)\in G_1'$ ; les inclusions instanciées en ces couples
donnent $(p,y)\in G_2$, $(y,r)\in G_2'$, d'où $(p,r)\in G_2'\circ G_2$
(axiome \texttt{COMPOSEE}, $\exists$-éliminations propres sur $p,r,y$).
Conditionnel honnête (2~hypothèses).

\textbf{Diagonale $\Delta_X$ (graphe de $\mathrm{Id}_X$, Déf.~8, E~II.13).} Deux résultats,
\emph{clos} (module \texttt{.../ii\_3\_correspondances/ensembles\_diagonale\_couple.py}) :
\begin{itemize}
  \item caractérisation au couple $\bigl((a,b)\in\Delta_X\bigr)\Leftrightarrow
        (a\in X\ \text{et}\ a=b)$ (\texttt{couple\_diagonale}) --- l'analogue, pour
        $\Delta$, de \texttt{couple\_reciproque}/\texttt{couple\_composee} ; preuve
        par l'axiome \texttt{DIAGONALE} et la Proposition~1 sur $(a,b)=(u,u)$ ;
  \item $\Delta_X^{-1}=\Delta_X$ (« $\mathrm{Id}_X$ est sa propre réciproque »,
        \texttt{diagonale\_auto\_reciproque}) --- \emph{inconditionnel} : $\Delta_X^{-1}$
        et $\Delta_X$ ne contiennent que des couples, double inclusion via les deux
        briques ci-dessus puis extensionnalité $A1$ ;
  \item $\mathrm{pr}_1\Delta_X=\mathrm{pr}_2\Delta_X=X$ (\texttt{pr1/pr2\_diagonale}) ---
        $\mathrm{dom}(\Delta_X)=\{z\mid(\exists y)(z\in X\wedge z=y)\}=X$ (effondrement
        du témoin), via la caractérisation au couple et \texttt{egalite\_par\_extension} ;
  \item composition avec l'identité, au couple :
        $\bigl((x,z)\in G\circ\Delta_A\bigr)\Leftrightarrow(x\in A\wedge(x,z)\in G)$ et le
        dual $\bigl((x,z)\in\Delta_B\circ G\bigr)\Leftrightarrow((x,z)\in G\wedge z\in B)$
        (\texttt{couple\_composee\_diagonale}, \texttt{diagonale\_composee\_couple}).
\end{itemize}

\textbf{Neutralité de l'identité (Déf.~8, E~II.13) : $\Gamma\circ\mathrm{Id}_A=
\mathrm{Id}_B\circ\Gamma=\Gamma$.} Au niveau des graphes, module
\texttt{.../ii\_3\_correspondances/ensembles\_identite\_neutre.py} :
\[
  \{\,\mathrm{est\_graphe}(G),\ \mathrm{pr}_1G\subset A\,\}\vdash G\circ\Delta_A=G,
  \qquad
  \{\,\mathrm{est\_graphe}(G),\ \mathrm{pr}_2G\subset B\,\}\vdash \Delta_B\circ G=G.
\]
\textbf{Preuve (\Nstar).} Aboutissement des briques au couple. Inclusion
$G\circ\Delta_A\subset G$ \emph{inconditionnelle} (l'étape $(p,y)\in\Delta_A$ force
$y=p$, d'où $(p,r)\in G$) ; $G\subset G\circ\Delta_A$ utilise $\mathrm{est\_graphe}(G)$
(tout $z\in G$ est un couple $(a,b)$) et $\mathrm{pr}_1G\subset A$ (pour $a\in A$),
puis recompose par \texttt{couple\_composee\_diagonale} ; extensionnalité $A1$.
Conditionnels honnêtes (les deux hypothèses sont load-bearing). $\mathrm{theorie}=22$.
